% Vietnamese translation for OI Public Library
% Source-File: IOI中国国家候选队论文/2005/钱自强/钱自强.doc
\documentclass[11pt]{article}
\usepackage{oipl-vn}
\usepackage{tikz}
\usetikzlibrary{positioning}

\newcommand{\GA}{\ensuremath{\mathrm{GA}}}
\newcommand{\Pruefer}{Prüfer}

\title{Phân tích và nghiên cứu ứng dụng của thuật toán di truyền}
\author{Qian Ziqiang\\Trường Trung học số 8 Fuzhou}
\date{2005}

\begin{document}
\maketitle

\origtitle{Phân tích và nghiên cứu ứng dụng của thuật toán di truyền}
\sourcepath{IOI China candidate team papers / 2005 / Qian Ziqiang.doc}

\renewcommand{\abstractname}{Tóm tắt}
\begin{abstract}
Cùng với sự phát triển của khoa học kỹ thuật, các bài toán gặp trong sản
xuất và đời sống ngày càng phức tạp. Nhiều bài toán đòi hỏi tìm nghiệm tối ưu
hoặc nghiệm xấp xỉ trong một không gian tìm kiếm rất lớn, nên việc dùng các
thuật toán truyền thống trở nên khó khăn. Những năm gần đây, tư tưởng tiến hóa
trong sinh học được ứng dụng rộng rãi trong kỹ thuật công trình, trí tuệ nhân
tạo và nhiều lĩnh vực khác, tạo thành một lớp thuật toán tìm kiếm ngẫu nhiên
hiệu quả gọi là thuật toán tiến hóa.

Bài viết trước hết giới thiệu nguyên lý và lịch sử của thuật toán tiến hóa,
sau đó trình bày chi tiết các bước dùng thuật toán di truyền để giải bài toán,
đồng thời nêu một số đề xuất cải tiến và kinh nghiệm áp dụng. Tiếp theo, bài
viết dùng một ví dụ tối ưu hóa NP-khó để minh họa sâu hơn cách dùng thuật toán
di truyền, rồi phân tích dữ liệu thử nghiệm để chứng minh hiệu quả của phương
pháp.
\end{abstract}

\noindent\textbf{Từ khóa:} trí tuệ nhân tạo; thuật toán tiến hóa; thuật toán
di truyền (\GA); cây khung nhỏ nhất đa mục tiêu.

\section{Lý thuyết thuật toán tiến hóa}

\subsection{Khái quát}

Từ thời xa xưa, sự sống bắt đầu từ các sinh vật đơn bào, trải qua các biến đổi
khắc nghiệt của môi trường và tiến hóa từ bậc thấp đến bậc cao, từ đơn giản
đến phức tạp. Sự sống không ngừng sinh sôi, cuối cùng tạo ra các sinh vật có
tư duy và trí tuệ. Con người có một cấu trúc sống ưu việt: không chỉ thụ động
thích nghi với môi trường, mà còn có thể học tập, bắt chước và sáng tạo để
liên tục nâng cao năng lực thích nghi.

Thuật toán tiến hóa là thuật toán tìm kiếm ngẫu nhiên mô phỏng chọn lọc tự
nhiên và cơ chế di truyền. Nó dùng quy luật ``ưu thắng liệt thải, thích nghi
thì tồn tại'' để khuyến khích các cấu trúc tốt, đồng thời dùng mô hình di
truyền và biến dị kiểu Mendel để giữ lại cấu trúc đã có trong quá trình lặp và
tìm kiếm cấu trúc tốt hơn. Là một phương pháp tối ưu hóa và tìm kiếm ngẫu
nhiên, thuật toán tiến hóa có các đặc điểm sau: nó không tìm kiếm mù quáng,
cũng không vét cạn toàn bộ không gian, mà dựa trên môi trường sống của cá thể,
tức hàm mục tiêu, để tìm kiếm có định hướng. Thuật toán chỉ cần giá trị của
mục tiêu, không cần thông tin khác, nên phù hợp với các bài toán tối ưu quy mô
lớn, phi tuyến mạnh, không liên tục và có nhiều cực trị cục bộ; tính phổ dụng
của nó rất cao. Đối tượng thao tác của thuật toán là một quần thể cá thể thay
vì một cá thể đơn lẻ, vì thế nó có nhiều quỹ đạo tìm kiếm cùng lúc.

\subsection{Thuật toán di truyền}

Thuật toán di truyền (\emph{Genetic Algorithm}) là một nhánh quan trọng của
thuật toán tiến hóa. Nó do John Holland đề xuất, ban đầu dùng để nghiên cứu
quá trình thích nghi của hệ tự nhiên và thiết kế phần mềm có tính tự thích
nghi. Gần đây, thuật toán di truyền đã được ứng dụng thành công như một công
cụ giải bài toán và tối ưu hóa, đồng thời ngày càng phổ biến.

Đặc điểm chính của thuật toán di truyền là chiến lược tìm kiếm theo quần thể
và việc trao đổi thông tin giữa các cá thể trong quần thể. Về bản chất, nó mô
phỏng quá trình học tập tổng thể của một quần thể gồm nhiều cá thể, trong đó
mỗi cá thể biểu diễn một điểm nghiệm trong không gian tìm kiếm. Từ một quần
thể ban đầu bất kỳ, thuật toán dùng các thao tác di truyền như chọn lọc, lai
ghép và đột biến để làm quần thể tiến hóa qua từng thế hệ tới những vùng ngày
càng tốt hơn của không gian tìm kiếm, cho đến khi đạt tới điểm nghiệm tối ưu.

So với các phương pháp tìm kiếm khác, ưu thế của thuật toán di truyền thể hiện
ở vài điểm. Trước hết, nó ít bị kẹt trong cực trị cục bộ; ngay cả khi hàm thích
nghi không liên tục hoặc không đều, thuật toán vẫn có xác suất lớn tìm được
nghiệm tối ưu toàn cục. Thứ hai, do có tính song song tự nhiên, thuật toán di
truyền rất phù hợp với xử lý song song và phân tán quy mô lớn. Ngoài ra, thuật
toán dễ kết hợp với các kỹ thuật khác như mạng nơ-ron, suy luận mờ, hành vi
hỗn độn và sự sống nhân tạo để tạo ra phương pháp giải bài toán có hiệu năng
tốt hơn.

\section{Thuật toán di truyền}

\subsection{Quy trình cơ bản}

Một thuật toán di truyền chạy tuần tự thường thực hiện các bước sau:
\begin{enumerate}
  \item Mã hóa bài toán cần giải; đặt \(t:=0\).
  \item Khởi tạo ngẫu nhiên quần thể
  \[
    X(0)=\{x_1,x_2,\ldots,x_M\}.
  \]
  \item Với mỗi nhiễm sắc thể \(x_i\) trong quần thể hiện tại \(X(t)\), tính
  độ thích nghi \(F(x_i)\). Độ thích nghi thể hiện năng lực thích ứng của cá
  thể với môi trường và quyết định xác suất cá thể ấy được rút chọn trong các
  thao tác di truyền.
  \item Áp dụng các toán tử di truyền lên \(X(t)\) theo xác suất đã đặt trước
  để tạo quần thể thế hệ mới \(X(t+1)\). Mục đích của các toán tử này là mở
  rộng phạm vi bao phủ của số cá thể hữu hạn, thể hiện tư tưởng tìm kiếm toàn
  cục.
  \item Đặt \(t:=t+1\), tức dùng thế hệ mới thay thế thế hệ cũ. Nếu chưa đạt
  điều kiện dừng đã định, quay lại bước 3.
\end{enumerate}

\subsection{Các yếu tố quan trọng}

\paragraph{Lý thuyết mã hóa.}
Thuật toán di truyền cần dùng một phương pháp mã hóa nào đó để ánh xạ không
gian nghiệm sang không gian mã, chẳng hạn chuỗi bit, số thực hoặc chuỗi có thứ
tự. Cách biểu diễn này tương tự cấu trúc nhiễm sắc thể sinh học, nên dễ giải
thích bằng lý thuyết di truyền và cũng dễ hiện thực các thao tác di truyền.
Lý thuyết mã hóa là một trong các yếu tố quyết định hiệu quả của thuật toán.
Mã hóa nhị phân là cách thường dùng nhất vì có nhiều mẫu toán tử và dễ cài
đặt. Tuy nhiên, trong từng bài toán cụ thể, nếu chọn được dạng mã phù hợp với
không gian nghiệm, ta có thể thao tác trực tiếp hơn trên kiểu hình của nghiệm,
dễ đưa thêm thông tin heuristic liên quan, và thường đạt hiệu quả cao hơn mã
hóa nhị phân.

\paragraph{Nhiễm sắc thể.}
Mỗi chuỗi mã tương ứng với một nghiệm cụ thể của bài toán, được gọi là một
nhiễm sắc thể hoặc một cá thể. Thông qua phương pháp mã hóa đã chọn, một nhiễm
sắc thể có thể tương ứng với một nghiệm cụ thể. Một số lượng nhiễm sắc thể cố
định tạo thành một thế hệ quần thể; các nhiễm sắc thể trong quần thể có thể
trùng nhau.

\paragraph{Khởi tạo ngẫu nhiên.}
Thế hệ đầu tiên có thể được sinh ngẫu nhiên hoặc sinh có quy luật, chẳng hạn
từ một điểm đã biết gần nghiệm rời rạc, hoặc bằng cách phân bố đều các nghiệm
khả thi. Trong một lần chạy thuật toán, việc chủ động khởi tạo quần thể nhiều
lần có thể giúp thuật toán có năng lực tìm kiếm nghiệm tối ưu toàn cục mạnh
hơn.

\paragraph{Độ thích nghi.}
Độ thích nghi của một nhiễm sắc thể là đánh giá năng lực sống sót của nó, và
quyết định xác suất nhiễm sắc thể đó được rút chọn. Hàm đánh giá là tiêu chuẩn
đo độ thích nghi. Các hàm thường gặp gồm: trực tiếp dùng trọng số của nghiệm,
chẳng hạn với bài toán ba lô \(0/1\), dùng tổng giá trị đồ vật được đưa vào ba
lô làm độ thích nghi; đếm số bit \(1\) hoặc \(0\) trong chuỗi nhị phân, với các
thuật toán dùng mã nhị phân; hoặc cộng điểm của nhiễm sắc thể trên từng mục
tiêu trong bài toán tối ưu đa mục tiêu. Với loại bài toán cuối, bài viết sẽ
đưa ra một hàm đánh giá hiệu quả hơn.

\paragraph{Toán tử di truyền.}
Toán tử di truyền là phần lõi của thuật toán. Chúng tác động trực tiếp lên
quần thể hiện có để sinh quần thể thế hệ kế tiếp, vì thế cách chọn và phối hợp
toán tử, cũng như tỷ lệ của từng toán tử, ảnh hưởng trực tiếp đến hiệu quả của
thuật toán. Một thuật toán di truyền thường gồm nhiều toán tử; mỗi toán tử
chạy độc lập, còn bản thân thuật toán chỉ quy định tỷ lệ đóng góp của từng
toán tử khi tạo thế hệ mới. Khi chạy, toán tử rút một số nhiễm sắc thể tương
ứng từ quần thể hiện tại, xử lý chúng, rồi đưa nhiễm sắc thể mới vào thế hệ
kế tiếp.

Một số toán tử di truyền thường gặp là:
\begin{itemize}
  \item \textbf{Toán tử giữ lại:} rút một nhiễm sắc thể và đưa thẳng vào thế
  hệ kế tiếp. Toán tử này giúp thuật toán hội tụ.
  \item \textbf{Toán tử lai ghép:} rút hai nhiễm sắc thể, trao đổi một số đoạn
  giữa chúng, rồi chọn cá thể có độ thích nghi cao hơn, hoặc chọn cả hai, để
  đưa vào thế hệ kế tiếp. Toán tử này giúp thuật toán tận dụng các nghiệm đã
  có để tìm nghiệm tốt hơn.
  \item \textbf{Toán tử đột biến:} rút một nhiễm sắc thể, thay đổi ngẫu nhiên
  một phần của nó, rồi đưa vào thế hệ kế tiếp. Toán tử này giúp thuật toán
  thoát cực trị cục bộ, tăng năng lực tìm kiếm toàn cục và tránh rơi vào tìm
  kiếm cục bộ. Xác suất của toán tử không nên quá cao, nếu không nghiệm sẽ bị
  phân tán và thuật toán không thể hội tụ.
\end{itemize}

\paragraph{Rút chọn.}
Rút chọn là một thao tác cơ bản quan trọng. Nó chọn các nhiễm sắc thể từ quần
thể hiện tại để cung cấp cho toán tử di truyền theo nguyên tắc ``ưu thắng liệt
thải'', dựa trên độ thích nghi của từng nhiễm sắc thể. Cách thường dùng là
vòng quay có thiên lệch. Giả sử quần thể có \(M\) cá thể, trước hết tính tổng
độ thích nghi
\[
  S=\sum_{i=1}^{M} F(x_i).
\]
Chia các số nguyên từ \(1\) đến \(S\) thành \(M\) đoạn, trong đó đoạn của nhiễm
sắc thể \(x_i\) có độ dài \(F(x_i)\). Khi sinh ngẫu nhiên một số nguyên từ
\(1\) đến \(S\), ta chọn nhiễm sắc thể ứng với đoạn chứa số đó. Như vậy xác
suất \(x_i\) được chọn là
\[
  p_i=\frac{F(x_i)}{S}.
\]

\paragraph{Điều kiện dừng.}
Điều kiện dừng dùng để tránh cho thuật toán lặp mãi. Thông thường có thể dừng
sau một số vòng lặp xác định, hoặc tự động dừng khi tốc độ hội tụ của nghiệm
thấp hơn một ngưỡng nào đó. Khi đạt điều kiện dừng, thuật toán kết thúc và trả
về nhiễm sắc thể tốt nhất thu được trong quá trình chạy, hoặc tất cả nghiệm
tối ưu không bị trội nếu bài toán yêu cầu.

\subsection{Thiết lập tham số quan trọng}

Khi dùng thuật toán di truyền, ta thường phải chọn tham số khác nhau cho từng
bài toán. Với các bài toán quy mô lớn, ngay trong một lần chạy cũng có thể đặt
các tham số khác nhau cho những giai đoạn khác nhau. Các tham số ảnh hưởng lớn
đến hiệu quả gồm:
\begin{itemize}
  \item \textbf{Kích thước quần thể:} số nhiễm sắc thể trong một thế hệ. Quần
  thể càng lớn thì chứa được càng nhiều loại nhiễm sắc thể, có lợi cho việc
  tìm kiếm nghiệm tối ưu toàn cục, nhưng tốc độ hội tụ giảm và thời gian cần
  thiết tăng.
  \item \textbf{Số vòng lặp:} số lần cập nhật quần thể tối đa. Tăng số vòng
  lặp giúp nghiệm hội tụ chính xác hơn nhưng tốn nhiều thời gian hơn. Nếu thời
  gian cho phép, chạy nhiều lần với khởi tạo lại quần thể thường hiệu quả hơn
  đặt số vòng lặp rất lớn trong một lần chạy.
  \item \textbf{Tỷ lệ giữ lại:} tỷ lệ của toán tử giữ lại, thường không vượt
  quá \(70\%\).
  \item \textbf{Tỷ lệ lai ghép:} tỷ lệ của toán tử lai ghép, thường không vượt
  quá \(50\%\).
  \item \textbf{Tỷ lệ đột biến:} tỷ lệ của toán tử đột biến, thường không vượt
  quá \(1\%\).
\end{itemize}

\section{Ứng dụng vào cây khung nhỏ nhất đa mục tiêu}

Nhiều bài toán đời sống như xây dựng đường sá, dựng mạng điện hay thiết kế
mạng đều có thể quy về mô hình cây khung nhỏ nhất. Các thuật toán kinh điển
như Prim và Kruskal đều giải được bài toán này với độ phức tạp tuyến tính theo
số cạnh nếu đã có cấu trúc dữ liệu phù hợp. Tuy nhiên, bài toán thực tế thường
không đơn giản như vậy. Một cạnh có thể mang nhiều trọng số: xây đường không
chỉ xét chiều dài mà còn phải xét môi trường và yếu tố nhân văn; dựng mạng
điện không chỉ xét chi phí dây mà còn xét độ khó thi công; nối mạng không chỉ
xét độ trễ mà còn xét độ ổn định và an toàn truyền tải. Khi đó bài toán trở
thành tìm các nghiệm tối ưu không bị trội của cây khung nhỏ nhất đa mục tiêu.
Bài toán này là NP-khó, các thuật toán thông thường như Prim và Kruskal không
còn đủ, còn tìm kiếm vét cạn thì quá đắt.

\subsection{Bài toán cây khung nhỏ nhất đa mục tiêu}

\subsubsection{Cây khung nhỏ nhất}

Trong lý thuyết đồ thị, một đồ thị con liên thông không có chu trình được gọi
là cây. Gọi \(T=(N,E_T)\) là một đồ thị với \(|N|\ge 3\). Sáu mô tả sau về cây
là tương đương:
\begin{enumerate}
  \item \(T\) liên thông và không có chu trình.
  \item \(T\) có \(|N|-1\) cạnh và không có chu trình.
  \item \(T\) liên thông và có \(|N|-1\) cạnh.
  \item \(T\) liên thông và mọi cạnh đều là cạnh cắt.
  \item Giữa hai đỉnh bất kỳ của \(T\) có đúng một đường đi.
  \item \(T\) không có chu trình, nhưng nếu thêm một cạnh giữa một cặp đỉnh
  không kề nhau thì tạo ra đúng một chu trình.
\end{enumerate}

Xét đồ thị vô hướng
\[
  G=(V,E,w),
\]
trong đó \(w\) là hàm trọng số. Nếu cây \(T=(V,E_T)\) chứa tất cả đỉnh của
\(G\), thì \(T\) là một cây khung của \(G\), có trọng số
\[
  W(T)=\sum_{e\in E_T} w(e).
\]
Một đồ thị có thể có nhiều cây khung; cây khung có tổng trọng số nhỏ nhất được
gọi là cây khung nhỏ nhất.

\subsubsection{Cây khung nhỏ nhất đa mục tiêu}

Khi mỗi cạnh trong đồ thị có nhiều trọng số, bài toán tương ứng được gọi là
bài toán cây khung nhỏ nhất đa mục tiêu. Mục tiêu của lớp bài toán này là tìm
tất cả cây khung tối ưu không bị trội. Ngay cả khi không thêm ràng buộc nào,
lớp bài toán này vẫn thuộc nhóm NP-khó.

Cho một đồ thị vô hướng liên thông \(G=(V,E)\), trong đó
\[
  V=\{v_1,v_2,\ldots,v_n\}
\]
là tập đỉnh hữu hạn. Với cạnh \(e_{ij}\in E\), giả sử cạnh này có \(m\) thuộc
tính dương
\[
  c(e_{ij})=(c^1_{ij},c^2_{ij},\ldots,c^m_{ij}),
\]
trong đó mỗi thuộc tính có thể là khoảng cách, chi phí, độ khó, v.v. Dùng biến
\(x_{ij}\in\{0,1\}\) để biểu diễn việc chọn cạnh \(e_{ij}\). Một vector \(x\)
biểu diễn một cây khung nếu các cạnh có \(x_{ij}=1\) tạo thành một cây trên
\(V\). Với mục tiêu thứ \(k\), chi phí của cây là
\[
  f_k(x)=\sum_{e_{ij}\in E} c^k_{ij}x_{ij}.
\]
Bài toán đa mục tiêu có thể viết là
\[
  \min_{x\in X} F(x)
  =
  \min_{x\in X} \bigl(f_1(x),f_2(x),\ldots,f_m(x)\bigr),
\]
trong đó \(X\) là tập tất cả cây khung khả thi. Ta cần tìm tập nghiệm không bị
trội: không tồn tại nghiệm khác tốt hơn hoặc bằng nó trên mọi mục tiêu và tốt
hơn hẳn trên ít nhất một mục tiêu.

So với cây khung nhỏ nhất thông thường, bài toán chỉ khác ở chỗ hàm mục tiêu
không còn đơn lẻ. Chính việc có nhiều mục tiêu, mà các mục tiêu này thường mâu
thuẫn với nhau, làm cho ta không thể dùng thuật toán kinh điển bằng cách lần
lượt chọn cạnh có trọng số nhỏ nhất. Nếu quy nhiều mục tiêu thành một mục tiêu
rồi áp dụng thuật toán kinh điển, ta chỉ thu được một nghiệm chứ không phải
một nhóm nghiệm tối ưu không bị trội; hơn nữa, bản thân việc chuyển đổi từ
nhiều mục tiêu sang một mục tiêu cũng là một khó khăn đối với người ra quyết
định.

\subsection{Dùng thuật toán di truyền để giải bài toán}

\subsubsection{Thiết kế mã hóa}

Cayley đã chứng minh rằng với một đồ thị đầy đủ \(G\), số cây nối tất cả \(n\)
đỉnh là \(n^{n-2}\). Dựa trên đó, \Pruefer{} xây dựng một song ánh giữa các
cây này và các chuỗi độ dài \(n-2\) lấy giá trị trong \(n\) số. Nếu đánh số
các đỉnh của đồ thị đầy đủ từ \(1\) đến \(n\), thì mỗi chuỗi độ dài \(n-2\)
trên tập \(\{1,2,\ldots,n\}\) tương ứng với đúng một cây khung.

Bài viết dùng mã \Pruefer{} cho cây khung: mỗi cây tương ứng với một chuỗi số
độ dài \(n-2\), và ngược lại mỗi chuỗi như vậy tương ứng với đúng một cây. Chứng
minh chi tiết của phép mã hóa từ cây sang chuỗi và phép giải mã từ chuỗi sang
cây có thể xem trong các tài liệu liên quan; ở đây chỉ mô tả ngắn gọn.

\paragraph{Quá trình mã hóa.}
\begin{enumerate}
  \item Ban đầu chuỗi mã rỗng.
  \item Gọi \(j\) là lá có số hiệu nhỏ nhất trong cây.
  \item Tìm đỉnh \(i\) duy nhất kề với \(j\), thêm \(i\) vào cuối chuỗi mã.
  \item Xóa \(j\) và cạnh nối \(i,j\) khỏi cây; lúc này cây còn \(n-1\) đỉnh.
  \item Lặp ba bước trên cho tới khi cây chỉ còn một cạnh. Chuỗi thu được là
  mã \Pruefer{} của cây.
\end{enumerate}

\paragraph{Quá trình giải mã.}
\begin{enumerate}
  \item Gọi \(P\) là chuỗi mã, và gọi \(S\) là tập các đỉnh của đồ thị không
  xuất hiện trong \(P\).
  \item Gọi \(i\) là đỉnh nhỏ nhất trong \(S\), còn \(j\) là phần tử đầu tiên
  bên trái của \(P\). Thêm cạnh \(ij\) vào cây; sau đó xóa \(i\) khỏi \(S\) và
  xóa phần tử \(j\) khỏi đầu \(P\). Nếu \(j\) không còn xuất hiện trong \(P\),
  thêm \(j\) vào \(S\).
  \item Lặp bước trên cho tới khi \(P\) rỗng.
  \item Khi \(P\) rỗng, trong \(S\) còn đúng hai đỉnh; thêm cạnh nối hai đỉnh
  này vào cây. Cây thu được chính là cây tương ứng với chuỗi \(P\) ban đầu.
\end{enumerate}

Hình dưới minh họa một cây khung và mã \Pruefer{} tương ứng. Khi mã hóa, ta
lần lượt xóa lá nhỏ nhất: \(1,2,3,4\), và ghi các đỉnh kề tương ứng
\(4,4,4,5\). Có thể thấy mã \Pruefer{} là một biểu diễn hiệu quả của cây
khung: nó cho phép sinh ngẫu nhiên một cây rất dễ, và thông tin ở từng vị trí
trong chuỗi mã tương đối đồng đều, nên rất phù hợp với các thao tác di truyền.

\begin{figure}[ht]
\centering
\begin{tikzpicture}[
  v/.style={circle, draw, minimum size=7mm, inner sep=0pt},
  node distance=1.1cm
]
\node[v] (4) {4};
\node[v, above left=of 4] (1) {1};
\node[v, left=of 4] (2) {2};
\node[v, below left=of 4] (3) {3};
\node[v, right=of 4] (5) {5};
\node[v, right=of 5] (6) {6};
\draw (1) -- (4) -- (2);
\draw (3) -- (4) -- (5) -- (6);
\node[below=0.8cm of 4, align=center]
  {Thứ tự xóa lá: \(1,2,3,4\)\\Mã \Pruefer: \((4,4,4,5)\)};
\end{tikzpicture}
\caption{Một cây khung sáu đỉnh và chuỗi \Pruefer{} thu được bằng quy tắc xóa lá nhỏ nhất.}
\end{figure}

\subsubsection{Hàm đánh giá}

Định nghĩa hàm đánh giá của nhiễm sắc thể \(x\) là
\[
  g(x)=\sum_{i=1}^{m}
  \left(\frac{\operatorname{min}[i]}{f_i(x)}\cdot 100\right)^2,
\]
trong đó \(f_i(x)\) là chi phí của nhiễm sắc thể hiện tại trên mục tiêu thứ
\(i\), còn \(\operatorname{min}[i]\) là chi phí nhỏ nhất trên mục tiêu thứ
\(i\) trong tất cả nhiễm sắc thể đã sinh ra cho tới trước thế hệ hiện tại.

Cách định nghĩa này tốt hơn việc cộng trực tiếp trọng số trên các mục tiêu ở
chỗ nó không chỉ phản ánh ưu thế của một nhiễm sắc thể trên từng mục tiêu, mà
còn tránh được hiện tượng ưu thế ở một mục tiêu bị che lấp do các mục tiêu có
miền giá trị hoặc xu hướng tổng thể khác nhau. Nhờ vậy thuật toán có thể thích
ứng tốt hơn với các loại bài toán thực tế.

\subsubsection{Định nghĩa thao tác di truyền}

\paragraph{Toán tử lai ghép.}
Toán tử lai ghép được dùng là lai ghép theo đoạn nhỏ cùng vị trí: chọn ngẫu
nhiên trong hai nhiễm sắc thể hai đoạn cùng vị trí, độ dài không quá \(2\), rồi
hoán đổi chúng và chọn cá thể có độ thích nghi cao hơn để đưa vào thế hệ kế
tiếp. Ví dụ tương ứng với hình 3-3 trong tài liệu gốc:
\[
\begin{array}{r|rrrrrr}
\text{Parent 1} & 2&5&6&8&2&5\\
\text{Parent 2} & 8&4&5&2&8&7\\ \hline
\text{Offspring 1} & 2&5&5&2&2&5\\
\text{Offspring 2} & 8&4&6&8&8&7
\end{array}
\]

\paragraph{Toán tử đột biến.}
Toán tử đột biến dùng đột biến một điểm thông thường: sinh ngẫu nhiên một số
từ \(1\) đến \(n\) để thay thế một vị trí nào đó trong chuỗi mã \Pruefer. Ví
dụ tương ứng với hình 3-4:
\[
\begin{array}{r|rrrrrr}
\text{Parent} & 2&5&6&8&4&5\\ \hline
\text{Offspring} & 2&5&6&2&4&5
\end{array}
\]
Có thể thấy đột biến một điểm vẫn giữ lại phần lớn tính chất của nhiễm sắc thể
ban đầu.

\subsection{Thử nghiệm}

Trong thử nghiệm, tác giả đặt tỷ lệ giữ lại là \(54\%\), tỷ lệ lai ghép là
\(45\%\), tỷ lệ đột biến là \(1\%\), đồng thời điều chỉnh số vòng lặp và kích
thước quần thể theo từng bộ dữ liệu. Thuật toán được so sánh với tìm kiếm
nguyên thủy. Để dễ quan sát mức xấp xỉ của thuật toán di truyền, mọi dữ liệu
đều dùng hai mục tiêu.

\noindent Môi trường thử nghiệm: P4 2.0GHz, 256 MB DDR, Windows XP, Delphi
7.0.

\subsubsection{Dữ liệu kinh điển quy mô nhỏ}

\begin{center}
\begingroup
\scriptsize
\setlength{\tabcolsep}{3pt}
\begin{tabular}{c|p{2.2cm}|p{1.8cm}|p{1.8cm}|p{2.2cm}|p{1.6cm}|p{1.8cm}}
\textbf{Mã} & \textbf{Thông tin} & \textbf{Tìm kiếm: thời gian} &
\textbf{Tìm kiếm: đúng} & \textbf{Tham số GA} &
\textbf{GA: thời gian} & \textbf{GA: đúng}\\ \hline
E1 & \(N=3,\ K=2,\ Ans=2\) & \(0\)s & hoàn toàn đúng &
\(L=20,\ P=10\) & \(0\)s & hoàn toàn đúng\\
E2 & \(N=4,\ K=2,\ Ans=6\) & \(0\)s & hoàn toàn đúng &
\(L=100,\ P=30\) & \(0\)s & hoàn toàn đúng\\
E3 & \(N=5,\ K=2,\ Ans=12\) & \(0\)s & hoàn toàn đúng &
\(L=1000,\ P=100\) & \(2\)s & hoàn toàn đúng
\end{tabular}
\endgroup
\end{center}

\subsubsection{Dữ liệu ngẫu nhiên quy mô vừa và lớn}

\begin{center}
\small
\begin{tabular}{c|p{4.5cm}|p{4.5cm}}
\textbf{Mã dữ liệu} & \textbf{H1} & \textbf{H2}\\ \hline
Quy mô & \(N=10,\ K=2\) & \(N=11,\ K=2\)\\
Tìm kiếm nguyên thủy & 22 phút & 6 giờ 20 phút\\
Tham số GA & \(L=20000,\ P=400\) & \(L=30000,\ P=400\)\\
Thời gian GA & 116s & 296s\\
Độ đúng & Thu được một nhóm nghiệm rất gần nghiệm tối ưu. &
Thu được một nhóm nghiệm khá gần nghiệm tối ưu.
\end{tabular}
\end{center}

\subsubsection{Phân tích dữ liệu đặc biệt}

Với tối ưu hóa hai mục tiêu, trong thực tế có thể tồn tại quan hệ giữa chi phí
của hai mục tiêu. Ví dụ, khi các thiết bị tương tự nhau, một liên kết mạng có
tốc độ tương đối cao thường có độ ổn định thấp hơn. Vì vậy tác giả thử nghiệm
hai trường hợp đặc biệt: tương quan thuận và tương quan nghịch.

\paragraph{Tương quan thuận.}
Nếu với hai cạnh bất kỳ \(e_a,e_b\) của đồ thị, điều kiện
\[
  c_1(e_a) < c_1(e_b)
\]
kéo theo
\[
  c_2(e_a) < c_2(e_b),
\]
thì trọng số của đồ thị được gọi là tương quan thuận.

\paragraph{Tương quan nghịch.}
Tương tự, nếu với hai cạnh bất kỳ \(e_a,e_b\), điều kiện
\[
  c_1(e_a) < c_1(e_b)
\]
kéo theo
\[
  c_2(e_a) > c_2(e_b),
\]
thì trọng số của đồ thị được gọi là tương quan nghịch.

\begin{center}
\small
\begin{tabular}{c|p{3.2cm}|p{3.2cm}|p{4.2cm}}
\textbf{Mã dữ liệu} & \textbf{S1} & \textbf{S2} & \textbf{S3}\\ \hline
Quy mô & \(N=10,\ K=2\) & \(N=10,\ K=2\) & \(N=15,\ K=2\)\\
Đặc tính & tương quan thuận & tương quan nghịch & tương quan thuận\\
Tìm kiếm nguyên thủy & 15 phút & 28 phút & khoảng 600 năm\\
Tham số GA & \(L=2000,\ P=400\) & \(L=20000,\ P=400\) &
\(L=20000,\ P=400\)\\
Thời gian GA & 11 giây & 304s & 114 giây mỗi lần, chạy 15 lần\\
Độ đúng & Trong 5 lần chạy có xác suất \(80\%\), tức 4 lần, thu được nghiệm
tối ưu. & Thu được một nhóm nghiệm khá tương tự nghiệm tốt, xem hình minh họa
trong tài liệu gốc. & Ở lần chạy thứ 12, chương trình thu được nghiệm tối ưu
\((3.11,3.83)\). Đáng chú ý là lần thứ 3 đã có nghiệm xấp xỉ rất gần
\((3.14,3.86)\), và ở lần thứ 6 cùng lần thứ 9 đều xuất hiện nghiệm khá tốt
\((3.18,3.99)\).
\end{tabular}
\end{center}

Các biểu đồ phân tích dữ liệu H1, H2, S1, S2 và S3 so sánh phân bố nghiệm;
nội dung định lượng chính của chúng được phản ánh trong các bảng trên.

\section{Kết luận}

Bài viết đã giới thiệu một số kiến thức cơ bản và kinh nghiệm sử dụng thuật
toán di truyền, đồng thời thông qua kết quả thử nghiệm cho thấy thuật toán di
truyền có ưu thế mạnh mà nhiều thuật toán khác khó sánh được khi giải các bài
toán tối ưu tổ hợp. Đặc điểm của nó là có thể thu được nghiệm tương đối thỏa
đáng trong thời gian ngắn, trong khi bản thân thuật toán khá đơn giản. Với rất
nhiều bài toán thực tế mà thuật toán thông thường không giải quyết được, thuật
toán di truyền có triển vọng ứng dụng tốt.

Thuật toán di truyền không chỉ là một thuật toán, mà còn là một cách tư duy.
Trong tìm kiếm, vận dụng linh hoạt tư tưởng tiến hóa thường đem lại hiệu quả
rất cao. Hiện nay thuật toán di truyền vẫn chưa thật phổ biến trong thi tin
học; mục đích của bài viết là giúp ngày càng nhiều người yêu thích tin học
hiểu thuật toán di truyền và tư tưởng của thuật toán tiến hóa.

Do năng lực của tác giả có hạn, bài viết khó tránh khỏi thiếu sót.

\section*{Tài liệu tham khảo}

\begin{enumerate}
  \item Chen Zhiping, Xu Zongben. \emph{Toán học máy tính}. Science Press.
  \item Zhou Ming, Sun Shudong. \emph{Nguyên lý và ứng dụng của thuật toán di
  truyền}. National Defense Industry Press.
  \item Shao Junli, Zhang Jing, Wei Changhua. \emph{Cơ sở trí tuệ nhân tạo}.
  Electronics Industry Press.
\end{enumerate}

\appendix
\section{Chương trình nguồn}

Phụ lục của tài liệu gốc đưa chương trình Pascal giải bài toán cây khung nhỏ
nhất đa mục tiêu bằng thuật toán di truyền. Mã dưới đây giữ nguyên cấu trúc
chính của chương trình, chỉ chuẩn hóa dấu nháy trong hai lệnh mở file để tránh
ký tự Word cũ.

\begin{verbatim}
Program Ga;
const
     MaxN          =    50;
     MaxK          =     5;
     PopulationMax =   1000;
     Maxans        =   1000;
     Population    =     40;
     Live          =   2000;
     Change        =     45;
     Suddenly      =      1;
type code=array[1..MaxN]of integer;
     vector=array[1..Maxk]of real;
     note=record
            v    :vector;
            data :code;
            power:longint;
          end;
     Group=array[1..PopulationMax]of note;
var
  inf,ouf:text;
  d:array[1..MaxK,0..MaxN,0..MaxN]of real;
  ans,ans2:array[1..Maxans]of vector;
  ansb,ansb2:array[1..maxans]of code;
  sn,n,k:longint; maxv:vector;
  population,live,change,suddenly:longint;

function Genetic_Uncode(s:code):vector;
var temp:vector; uu:boolean;
    i,j,l,r:longint;
    s2:code; Se:set of 1..maxn;
  function get:integer;
  var i:longint;
  begin
    for i:=1 to n do
      if i in Se then begin Se:=Se-[i]; get:=i; exit; end;
  end;
begin
  Se:=[1..n];
  for i:=1 to k do temp[i]:=0;
  for i:=1 to n-2 do Se:=Se-[s[i]];
  for i:=1 to n-2 do
  begin
    l:=get;
    for j:=1 to k do
    begin
      if (l<=0)or(l>n) or (s[i]<=0)or(s[i]>n) then
      begin
        L:=L+1;
      end;
      temp[j]:=temp[j]+d[j,l,s[i]];
    end;
    uu:=true;
    for j:=i+1 to n-2 do
      if s[j]=s[i] then begin uu:=false; break; end;
    if uu then Se:=Se+[s[i]];
  end;
  l:=get; r:=get;
  for j:=1 to k do
    temp[j]:=temp[j]+d[j,l,r];
  Genetic_Uncode:=temp;
end;

procedure Input_Initial;
var s,i,j:longint; tempcode:code;
begin
  readln(inf,n,k);
  for s:=1 to k do
    for i:=1 to n do
      for j:=1 to n do
        read(inf,d[s,i,j]);
  for i:=1 to n-2 do tempcode[i]:=random(n)+1;
  maxv:=Genetic_uncode(tempcode); randomize;
  sn:=1; ans[1]:=maxv; ansb[1]:=tempcode;
  close(inf);
end;

function find(v1,v2:vector):longint;
var i:longint;
    check:boolean;
begin
  check:=true;
  for i:=1 to k do
    if v1[i]<v2[i] then begin check:=false; break; end;
  if check then begin find:=1; exit; end;
  check:=true;
  for i:=1 to k do
    if v1[i]>v2[i] then begin check:=false; break; end;
  if check then begin find:=-1; exit; end;
  find:=0;
end;

procedure insert(v:vector; vcode:code);
var i,sn2:longint;
begin
  sn2:=0;
  if (v[1]=2) then begin end;
  for i:=1 to sn do
    case find(v,ans[i]) of
      0: begin inc(sn2); ans2[sn2]:=ans[i]; ansb2[sn2]:=ansb[i]; end;
      1: exit;
    end;
  inc(sn2); ans2[sn2]:=v; ansb2[sn2]:=vcode;
  sn:=sn2; ans:=ans2; ansb:=ansb2;
end;

procedure Genetic;
var S,S0:Group;
    i:longint; all:int64;

  procedure Genetic_Initial;
  var i,j:longint;
  begin
    for i:=1 to Population do
      for j:=1 to n-2 do
        s[i].data[j]:=random(n)+1;
  end;

  function Genetic_Compute(v:vector; vcode:code):longint;
  var i,mark:longint; nowv:vector;
  begin
    nowv:=maxv; mark:=0; insert(v,vcode);
    for i:=1 to k do
    begin
      mark:=mark+sqr(round((nowv[i]/v[i])*100));
      if v[i]<maxv[i] then maxv[i]:=v[i];
    end;
    Genetic_Compute:=mark;
  end;

  procedure Genetic_Start;
  var i:longint;
  begin
    all:=0;
    for i:=1 to Population do
    begin
      s[i].v:=Genetic_Uncode(s[i].data);
      s[i].power:=Genetic_Compute(s[i].v,s[i].data);
      all:=all+s[i].power;
    end;
  end;

  function Genetic_Random:code;
  var now:longint; seed:int64;
  begin
    Seed:=random(all)+1; now:=1;
    while Seed>s[now].power do
    begin
      Seed:=Seed-s[now].power; inc(now);
    end;
    Genetic_Random:=s[now].data;
  end;

  function Genetic_Suddenly:code;
  var temp:code;
      Seed:longint;
  begin
    temp:=Genetic_Random; Seed:=random(N-2)+1;
    Temp[Seed]:=random(n)+1;
    Genetic_Suddenly:=temp;
  end;

  function Genetic_Change:code;
  var temp1,temp2:code;
      Seed1,Seed2,temp:longint;
  begin
    temp1:=Genetic_Random; Seed1:=random(N-2)+1;
    temp2:=Genetic_Random; Seed2:=random(N-2)+1;
    Temp:=temp1[Seed1];
    Temp1[Seed1]:=Temp2[Seed2];
    Temp2[Seed2]:=Temp;
    if Genetic_Compute(Genetic_Uncode(temp1),temp1)>
       Genetic_Compute(Genetic_Uncode(temp2),temp2) then
      Genetic_Change:=temp1
    else
      Genetic_Change:=temp2;
  end;

  procedure Genetic_Genarate;
  var Seed,i:longint;
  begin
    Genetic_Start;
    for i:=1 to Population do
    begin
      Seed:=random(100)+1;
      if Seed <= Suddenly then
        s0[i].data:=Genetic_Suddenly
      else if seed<= Change+Suddenly then
        s0[i].data:=Genetic_Change
      else
        s0[i].data:=Genetic_Random;
    end;
    s:=s0;
  end;

begin
  Genetic_Initial;
  for i:=1 to Live do
  begin
    Genetic_Genarate;
  end;
end;

procedure sortans(l,r:longint);
var tl,tr:longint; tmp:vector;
begin
  tl:=l; tr:=r; tmp:=ans[l];
  while tl<tr do
  begin
    while (ans[tr][1]>tmp[1])and(tr>tl) do dec(tr);
    if tr=tl then break; ans[tl]:=ans[tr]; inc(tl);
    while (ans[tl][1]<tmp[1])and(tl<tr) do inc(tl);
    if tr=tl then break; ans[tr]:=ans[tl]; dec(tr);
  end;
  ans[tl]:=tmp;
  if tl>l+1 then sortans(l,tl-1);
  if tr+1<r then sortans(tr+1,r);
end;

procedure Output_Result;
var i,j:longint;
begin
  sortans(1,sn);
  for i:=1 to sn do
  begin
    write(ouf,'The ',i,'th Answer:');
    for j:=1 to k do
      write(ouf,ans[i][j]:9:2);
    write(ouf,'  Code=');
    for j:=1 to n-2 do write(ouf,' ',ansb[i][j]);
    writeln(ouf);
  end;
  writeln(ouf);
  for j:=1 to k do
  begin
    for i:=1 to sn do
      writeln(ouf,ans[i][j]:9:2);
    writeln(ouf,'------');
  end;
  writeln(ouf,'Total=',sn);
  close(ouf);
end;

procedure start;
begin
  assign(inf,'input.txt'); reset(inf);
  assign(ouf,'edit8.text'); rewrite(ouf);
  Input_Initial;
  Genetic;
  Output_Result;
end;

begin
  start;
end.
\end{verbatim}

\end{document}
